\input{_preamble}
\begin{document}
% The shape a real paper actually writes: the free translation is a third
% \\-separated line after \gll, with no \glt.  linguexx typesets that
% identically to \glt -- measured, same x and same baseline -- so it is the
% translation and not text that precedes the example.
\ex. \gll quodsi magnam Hermagoras habuisset facultatem \\
     if great.\textsc{acc} Hermagoras.\textsc{nom} have.\textsc{pluprf} ability.\textsc{acc} \\
     `if Hermagoras had acquired great skill' (Cicero)

% Same for \glll: a fourth line is the translation, not a fourth tier.
\ex. \glll Ich habe geschlafen \\
      ich hab-e schlaf-en \\
      \lpzg{1sg} have-\lpzg{1sg} sleep-\lpzg{ptcp} \\
      `I slept.'

% Genuine text BEFORE the gloss still is that, and still warns -- the
% trailing line must not swallow it.
\ex. As Cicero puts it, \gll magnam facultatem \\
     great.\textsc{acc} ability.\textsc{acc} \\
     `great skill'

% A judgment mark in front of the gloss, with a trailing translation: the
% mark must reach the judgment slot and must not drag the translation into
% tier 0 behind it.
\ex. *\gll magnam facultatem \\
     great.\textsc{acc} ability.\textsc{acc} \\
     `great skill'

% Both a trailing line and a \glt: linguexx prints the trailing text and
% then the \glt line, in that order.  One slot, so they are joined, and
% that loses a line break -- which is worth a warning.
\ex. \gll magnam facultatem \\
     great.\textsc{acc} ability.\textsc{acc} \\
     `great skill'
\glt (Cicero, \emph{De inventione})
\end{document}
